%=============================================================================!
% 14.7  A PISTON WITH MASS                                                    !
%=============================================================================!
\section{A piston with mass}
\label{sec:pr-piston}

The barostats of \cref{sec:pr-rescaling} move the volume as a heavily
damped quantity, with no momentum of its own. A real piston has mass: it
overshoots and swings. This section builds the barostat that treats the
volume as a coordinate with inertia, finds the frequency at which it
swings, and measures it.

\subsection*{A piston on a gas spring}

A piston of mass $M_{\mathrm p}$ and area $A$ rests on a column of gas of
height $x$, so $V = Ax$, with the pressure $P_0$ above it. The gas pushes
up with $PA$ and the outside pushes down with $P_0A$, so $M_{\mathrm
p}\ddot x = A(P - P_0)$. If the piston is moved by $\delta x$ from where
the two balance, the volume changes by $A\,\delta x$ and the gas's
pressure by $\delta P = -A\,\delta x/(V\kappa)$, and
\begin{equation*}
   M_{\mathrm p}\,\ddot{\delta x} = -\frac{A^2}{V\kappa}\,\delta x :
\end{equation*}
the oscillator of \cref{sec:os-small}, with the gas as its spring, of
stiffness $A^2/V\kappa$, and the frequency $\omega^2 = A^2/(M_{\mathrm
p}V\kappa)$. A heavy piston on a soft gas swings slowly.

\subsection*{The volume as a coordinate}

Andersen turned this into a barostat\cite{Andersen1980}: the volume is
given a mass and a kinetic energy of its own, the outside does the work
$P_0V$, and the atoms' positions scale with the box, so that the volume
accelerates in proportion to $P - P_0$ and the sum of $H$, $P_0V$ and the
piston's kinetic energy is kept. The form of Martyna, Tobias and Klein
(MTK)\cite{Martyna1994}, which ASE and \texttt{mdlab} follow, takes as its coordinate the strain $\epsilon$
by which every length is $\mathrm{e}^\epsilon$ times its starting value, so
$V = V_0\mathrm{e}^{3\epsilon}$, with the momentum $p_\epsilon$ and the
mass $M_\epsilon$. Its equations are
\begin{align*}
   \dot{\vb{r}}_i &= \vb{v}_i + \frac{p_\epsilon}{M_\epsilon}\vb{r}_i ,
   & m_i\dot{\vb{v}}_i &= \vb{F}_i - \Bigl(1 + \frac{3}{N_{\mathrm f}}\Bigr)
   \frac{p_\epsilon}{M_\epsilon}m_i\vb{v}_i ,\\
   \dot V &= \frac{3Vp_\epsilon}{M_\epsilon} ,
   & \dot p_\epsilon &= 3V(P - P_0) + \frac{3}{N_{\mathrm f}}\sum_im_iv_i^2 ,
\end{align*}
with $N_{\mathrm f}$ the number of free momenta. Here $\vb{v}_i =
\vb{p}_i/m_i$ is the velocity of atom $i$ relative to the stretching box,
which the first equation adds to the box's own motion to give
$\dot{\vb{r}}_i$, and $P$ is the instantaneous pressure computed with
these $\vb{v}_i$. The first equation carries each atom along as the box
grows, the second slows the atoms as it does, the third ties the volume
to the strain, and the fourth pushes the piston with the difference of
pressures. Martyna, Tobias and Klein showed that with the two terms in
$3/N_{\mathrm f}$ these equations, given a thermostat, sample the ensemble
at constant pressure; without them the distribution of the volume is
wrong by an amount that shrinks as $N$ grows. Nosé-Hoover chains
(\cref{sec:th-nose}) add a friction to the second and fourth equations,
acting on the atoms and on $p_\epsilon$, and the run keeps $K + U + P_0V +
p_\epsilon^2/2M_\epsilon$ plus the chains' energies.
\texttt{barostats.run\_mtk} integrates these equations by the Trotter
splitting (\cref{sec:in-splitting}) of Tuckerman and
colleagues\cite{Tuckerman2006}, as ASE's
\texttt{IsotropicMTKNPT} does, and follows it to \SI{5.6e-7}{\angstrom}
over 40 steps. Both take $N_{\mathrm f} = 3N$. A liquid whose total
momentum stays at zero has only $3N - 3$ free momenta, and with $3N$ in
the chain its kinetic temperature runs high by about $3/(3N - 3)$, 0.4\%
for 256 atoms: the three runs of \cref{fig:pr-volume}(b) average 135.48,
135.99 and \SI{135.94}{\kelvin} against the \SI{135}{\kelvin} asked
for.

\subsection*{How fast it swings}

Ignoring the chains and the term $(3/N_{\mathrm f})\sum_im_iv_i^2$, which
stays close to $3\kB T$ and only shifts where the piston balances,
$M_\epsilon\ddot\epsilon = 3V(P - P_0)$. A change of the strain by
$\delta\epsilon$ changes the volume by $3V\delta\epsilon$ and, taking the
response at fixed temperature, the pressure by $\delta P =
-3\delta\epsilon/\kappa_T$, so
\begin{equation*}
   M_\epsilon\,\ddot{\delta\epsilon} = -\frac{9V}{\kappa_T}\,\delta\epsilon ,
   \qquad \omega^2 = \frac{9V}{M_\epsilon\kappa_T} ,
\end{equation*}
the piston on its gas spring again. ASE and \texttt{mdlab} take the mass
from a time constant, $M_\epsilon = (3N + 3)\kB T\tau_{\mathrm P}^2$, which
gives the period (\cref{ex:pr-period})
\begin{equation*}
   \frac{2\pi}{\omega} = 2\pi\tau_{\mathrm P}\sqrt{\frac{(N + 1)\kB T
   \kappa_T}{3V}} ;
\end{equation*}
for the liquid, with $\kappa_T = \SI{1.46}{\per\giga\pascal}$, it is
$0.865\,\tau_{\mathrm P}$, and $M_\epsilon = \SI{8.97e6}{\electronvolt\femto
\second\squared}$ for $\tau_{\mathrm P} = \SI{1}{\pico\second}$. Measured
from the crossings of the volume through its mean
(\cref{fig:pr-piston}), the periods are 0.248, 0.876, 2.450 and
\SI{8.695}{\pico\second} for $\tau_{\mathrm P}$ of 0.3, 1, 3 and
\SI{10}{\pico\second}, against 0.259, 0.865, 2.594 and
\SI{8.647}{\pico\second} predicted, the agreement within 6\% bearing out
the response taken at fixed temperature.

\begin{figure}[htb]
\centering
\includegraphics{ch14_pressure/piston}
\caption{A piston with mass rings. The liquid at \SI{135}{\kelvin} and
\SI{0.1}{\giga\pascal} under the MTK piston with chains. (a) The volume's
departure from its mean over the first \SI{30}{\pico\second} for
$\tau_{\mathrm P} = 1$, 3 and \SI{10}{\pico\second}, each trace offset by
\SI{1200}{\angstrom\cubed}. (b) The period of the volume's swing against
$\tau_{\mathrm P}$ (points) and the period of the linearised piston
(dashed).}
\label{fig:pr-piston}
\end{figure}

The piston samples the right volume (\cref{fig:pr-volume}): on the ideal
gas its mean is 0.999 and its spread 0.994 of the exact values, and on
the liquid, with $\tau_{\mathrm P} = \SI{1}{\pico\second}$, the volume's
fluctuations give $\kappa_T = (1.39 \pm 0.06)$\,\si{\per\giga\pascal}
against \SI{1.46}{\per\giga\pascal}. How heavy the piston should be is a
question of time. With $\tau_{\mathrm P} = \SI{10}{\pico\second}$ the
volume swings only 10 times in \SI{90}{\pico\second}, too few for a run of
that length to average its swings away; with \SI{1}{\pico\second} it
swings 103 times.

\subsection*{Changing shape}

Parrinello and Rahman gave every component of the cell matrix $\vb{h}$ a
mass and a momentum, so that the cell can change its shape as well as its
size, each lattice vector a piston of its own, and used it to follow a
crystal from one structure to another\cite{Parrinello1981}. Martyna,
Tobias and Klein gave the same freedom to their piston. \texttt{mdlab}
has only the isotropic piston; \cref{sec:pr-shape} changes the shape of
the cell with first-order barostats instead.

\begin{keyidea}
A piston with mass treats the volume as a coordinate pushed by $P -
P_0$; it swings like a mass on a gas spring, with $\omega^2 = 9V/(M_\epsilon
\kappa_T)$, a period of $0.865\,\tau_{\mathrm P}$ for the liquid with the
mass $(3N + 3)\kB T\tau_{\mathrm P}^2$. With Nosé-Hoover chains and the
corrections of Martyna, Tobias and Klein it samples the right volume; a
heavy piston swings so slowly that a short run averages few swings.
\end{keyidea}

\notebook{14}{set the piston's mass and watch the volume ring}

\begin{selfcheck}
A piston of \SI{1}{\kilo\gram} and area \SI{10}{\centi\metre\squared}, its
weight ignored, sits on one litre of air, taken at a fixed temperature, so that $\kappa
= 1/P$ with $P = \SI{1}{bar}$. What is the period of its swing?
\end{selfcheck}
